\section{Curved motion}\label{sec:motion}

This section constrains the head motion of every linear-time transposer, oblivious or adaptive; the decomposition of
the oblivious case follows in \cref{sec:decomposition}. The first paper~\cite{FirstPaper2026} rescales the work-head trajectories of a machine, passes to uniform limits
along subsequences of sizes, and proves that the limits of a linear-time transposer satisfy an isoperimetric
inequality $(\star)$. This is Theorem~M*, restated as \cref{thm:Mstar}. In this section we turn $(\star)$ into a
constraint on the shape of the motion. We combine it with the second paper's labels for finitely many affine pieces or
affine maps, built from F\o lner sets of finitely many affine maps and Lean-verified there~\cite{SecondPaper2026}, and
with the area formula, which charges the remaining returns of the heads by their speed (\cref{sec:mot-tools}). The outcome is Theorem~Cq (\cref{thm:Cq}): in every limit of a linear-time transposer, the work heads travel a total
distance of at least $\frac12$, in block units, along motion that is not countably-piecewise affine. Theorem~C
(\cref{thm:C}) applies this to machines whose work heads are all home-based, with legs made of boundedly many runs at
constant rate: no such machine, oblivious or adaptive, transposes in linear time.
\Cref{sec:mot-consequences} draws consequences. One work tape never suffices. Theorem~Cr (\cref{thm:Cr}) passes from
affine pieces to affine returns: in every common clock, that is, every reparametrization of time by a homeomorphism
whose inverse maps null sets to null sets, and for every countable set of affine maps, the times at which a head's next
return to its level is not given by one of these maps carry at least one block unit of travel in total; so a common
change of clock and fixed changes of scale do not help. Machines whose work heads all follow stroke trees are excluded, and a construction
must use head-dependent curved motion on at least two interacting tapes (\cref{cor:mot-design}).

These are restricted impossibility results, and they are unconditional: they use neither depth reduction
(\cref{prob:tcp}) nor any statement about codes. On the construction side they say which motion a linear-time
transposer must have. No numbered result of this section is formalized in this paper's Lean library. Theorem~M* and the finite affine
labels, on which they rest, are Lean-verified in the first two papers; the remaining cited ingredient, the area formula
for Lipschitz functions, is classical, as are the other facts of measure theory used below. Amenability enters only
through F\o lner sets of finitely many affine maps, so no hyperfiniteness theorem is needed.

\subsection{Rescaled limits and Theorem M*}\label{sec:mot-setting}

In this section the letter $t$ indexes the work heads, from $0$ as in the first paper, and $h$ denotes a home cell.
Let $M$ be a machine with $k\ge1$ work tapes whose work alphabet has $g+1$ symbols, and let $G=\max(1,\clog{(g+1)})$
be the number of bits per work cell~\cite[Definition~6.7]{FirstPaper2026}. For $d\ge1$ put $r_d=\lceil d\,b(d)/G\rceil$
and $R_d=d\,r_d$, so that $L(d)/G\le R_d\le L(d)/G+d$. Time is cut into \emph{slots} of $r_d$ steps and each work tape
into \emph{blocks} of $r_d$ cells; a block holds about one row of the matrix. On an input $w$ of size $d$ let
$x_{t,w}(s)$ be the cell of work head~$t$ after $s$ steps, frozen once $M$ halts. The \emph{rescaled trajectory} of
head~$t$ is
\[
  P^{(d)}_{t,w}(u)=\frac1d\Bigl\lfloor\frac{x_{t,w}(\lfloor ud\rfloor r_d)}{r_d}\Bigr\rfloor\qquad(u\ge0),
\]
the block of the head at the first step of slot $\lfloor ud\rfloor$, in units of $d$ blocks~\cite[Section~7]{FirstPaper2026}.
Macroscopic time and space are thus both measured in units of $R_d$ steps or cells. We fix a time constant $c>0$ and
put $c_0=cG+1$; a run of at most $cL(d)$ steps ends by macroscopic time $cL(d)/R_d\le cG<c_0$.

A \emph{trace class} of $M$ at size $d$~\cite[Definition~6.1]{FirstPaper2026} is a set $F$ of matrix inputs of size $d$
with a \emph{representative} $w\in F$ such that all members of $F$ halt after the same number of steps and have the
same block-level trace as $w$: slot by slot, the same block residencies of the work heads, the same input rows touched
and the same output length. Its \emph{co-entropy} is the least $D$ with $2^{L(d)}\le2^D\card F$. If $M$ is oblivious,
all matrix inputs of size $d$ form one trace class, of co-entropy~$0$~\cite[Lemma~6.2]{FirstPaper2026}.

\begin{definition}[Limits of a machine]\label{def:mot-limit}
A family $(P_t)_{t<k}$ of functions $[0,c_0]\to\R$ is a \emph{limit of $M$ at time constant $c$} if there are sizes
$d_i\to\infty$ and, for all large $i$, trace classes $F_i$ of $M$ at size $d_i$ with representatives $w_i$, common
halting times $T_i\le cL(d_i)$ and co-entropies $D_i$ with $D_i/L(d_i)\to0$, such that $M$ outputs the transpose on
every input in $F_i$ and $P^{(d_i)}_{t,w_i}\to P_t$ uniformly on $[0,c_0]$ for every $t<k$.
\end{definition}

For a limit, and more generally for any family of $1$-Lipschitz functions $P_t$ on $[0,c_0]$, we write $\lambda$ for
Lebesgue measure and
\[
  \nu_t=\abs{P_t'(u)}\,\mathrm du\ \text{ on }[0,c_0],\qquad
  \Lambda_t(p)=\{u\in[0,c_0]:P_t(u)=p\}.
\]
The derivative $P_t'$ exists almost everywhere and $\abs{P_t'}\le1$, so $\nu_t\le\lambda$. We call $\nu_t$ the
\emph{speed measure} of head~$t$: $\nu_t(A)$ is the distance travelled by the limit path during the times in $A$.
Distance $1$, which we call a \emph{block unit}, is $d$ blocks, that is, $R_d\approx L(d)/G$ cells.

The inequality $(\star)$ of the first paper uses labellings by finitely many intervals. An \emph{excised set} is a
finite union $S$ of closed intervals. A \emph{labelling} is a map $\beta\colon\R\to\N$ that is \emph{piecewise constant
off $S$}: for some finite set $Z$, $\beta$ is constant on every interval $[u,v]\subseteq[0,c_0]\setminus(S\cup Z)$. Its
\emph{classes off $S$} are the sets $\{u\in[0,c_0]\setminus S:\beta(u)=\ell\}$, $\ell\in\N$. Two points $s<s'$ of
$\Lambda_t(p)$ are \emph{consecutive} if no point of $\Lambda_t(p)$ lies strictly between them. Let $\kappa_t(p)$ be the
number of consecutive pairs $s<s'$ of $\Lambda_t(p)$ with $s,s'\notin S$ and $\beta(s)\ne\beta(s')$ if $\Lambda_t(p)$ is
finite, and $\kappa_t(p)=0$ if $\Lambda_t(p)$ is infinite. For $\theta\in[0,\infty]$ the inequality reads
\[
  (\star)\qquad 1\;\le\;\sum_{t<k}\int\kappa_t(p)\,\mathrm dp\;+\;2\theta\;+\;2k\,\lambda(S\cap[0,c_0]),
\]
where $\int$ denotes the lower Lebesgue integral, which is defined and monotone for every function with values in
$[0,\infty]$, measurable or not. We say that $(P_t)_{t<k}$ \emph{satisfies $(\star)_\emptyset$} if $(\star)$ holds with
$S=\emptyset$ for every $\theta\in(0,\frac12)$ and every labelling whose classes have measure at most $\theta$.

\begin{theorem}[Theorem M*~\cite{FirstPaper2026}]\label{thm:Mstar}
\cited{proved and Lean-verified in~\cite{FirstPaper2026}; not reproved here.}
Let $M$ be a machine with $k\ge1$ work tapes and let $c>0$.
\begin{enumerate}[label=(\alph*)]
\item Every limit $(P_t)_{t<k}$ of $M$ at time constant $c$ consists of $1$-Lipschitz functions on $[0,c_0]$ and
satisfies $(\star)$ for every $\theta\in[0,\infty]$, every excised set $S$ and every labelling, piecewise
constant off $S$, whose classes off $S$ have measure at most $\theta$. In particular it satisfies $(\star)_\emptyset$.
\item Let $d_i\to\infty$, and suppose that for all large $i$ the machine $M$ transposes every matrix of size $d_i$ and
halts within $cL(d_i)$ steps on it. Then $M$ has a limit at time constant $c$ along a subsequence of $(d_i)$, with trace
classes of co-entropy at most $C\,d_i\,b(d_i)$ for a constant $C$ that depends only on $M$ and $c$.
\end{enumerate}
\end{theorem}

This combines Lemma~7.1(a), Theorems~7.8 and~7.9 and the popular trace theorem (Theorem~6.8) of~\cite{FirstPaper2026};
the Lean declarations there are \leanname{Transpose.Tape.theoremMStar_level_class_on} for~(a) (with
\leanname{Transpose.Tape.blockTraj_limit_lipschitz} for the Lipschitz property) and
\leanname{Transpose.Tape.theoremMStar_level_adaptive} for~(b). The classes in~(b) are the \emph{popular} trace classes;
for an oblivious machine one can take all matrix inputs. A machine that witnesses~(T) with constant $c$ transposes
within $2cL(d)$ steps on every matrix input of every size, because $c(L+1)\le2cL$ (see the proof of Corollary~7.11
of~\cite{FirstPaper2026}; Lean: \leanname{Transpose.Tape.adaptive_hyps_of_linearTime}), so by~(b) it has limits at time
constant $2c$ along a subsequence of every sequence of sizes. We use only the case $S=\emptyset$ of~(a).

\subsection{Levels, returns and finite affine labels}\label{sec:mot-tools}

Statements about $1$-Lipschitz functions, here and below, hold for every $c_0>0$; for walks and limits of machines,
$c_0=cG+1$. In this subsection $P$ is a $1$-Lipschitz function on $[0,c_0]$, $\nu=\abs{P'}\,\mathrm du$ and
$\Lambda(p)=\{u\in[0,c_0]:P(u)=p\}$; for a family $(P_t)$ we add the index $t$ to everything.

\begin{lemma}[Levels]\label{lem:mot-levels}
\notformalized
\begin{enumerate}[label=(\alph*)]
\item For every Borel set $A\subseteq[0,c_0]$ the function $p\mapsto\#(A\cap\Lambda(p))\in\{0,1,\dots,\infty\}$ is
Lebesgue measurable, and $\int_\R\#(A\cap\Lambda(p))\,\mathrm dp=\nu(A)$.
\item $\Lambda(p)$ is finite for almost every $p$, and $\nu(\Lambda(p))=0$ for every $p$.
\item $\nu(A)\le\lambda(A)$ for every Borel set $A\subseteq[0,c_0]$; in particular $\lambda$-null sets are $\nu$-null.
\end{enumerate}
\end{lemma}

\begin{informalproof}
(a) Extend $P$ to a $1$-Lipschitz function on $\R$ that is constant on each side of $[0,c_0]$. The area formula for
Lipschitz maps from $\R$ to $\R$~\cite[Theorem~3.2.3]{Federer1969} states that for every Lebesgue measurable set
$A\subseteq\R$ the multiplicity function $p\mapsto\#(A\cap P^{-1}(p))$ is measurable and has integral $\int_A\abs{P'}$.
For $A\subseteq[0,c_0]$ the multiplicity is $\#(A\cap\Lambda(p))$. (For intervals $A$ this is Banach's indicatrix
theorem~\cite{Banach1925}.)
(b) For $A=[0,c_0]$, (a) gives $\int\#\Lambda(p)\,\mathrm dp=\nu([0,c_0])\le c_0<\infty$, so $\#\Lambda(p)<\infty$ for
almost every $p$. For $A=\Lambda(p_0)$, which is closed, the integrand in (a) vanishes for $p\ne p_0$, so
$\nu(\Lambda(p_0))=0$.
(c) $\abs{P'}\le1$ almost everywhere.
\end{informalproof}

All labellings below are piecewise constant off $S=\emptyset$, and $\kappa$ is the crossing count of $(\star)$ with
$S=\emptyset$. For a set $R\subseteq[0,c_0]$ of \emph{retained} times, let $\kappa^R(p)$ be the number of consecutive
pairs $s<s'$ of $\Lambda(p)$ with $s,s'\in R$ and $\beta(s)\ne\beta(s')$ if $\Lambda(p)$ is finite, and $\kappa^R(p)=0$
otherwise; so $\kappa=\kappa^{[0,c_0]}$. The points of $\Lambda(p)$ outside $R$ are the \emph{omitted} visits.

\begin{lemma}[Deleting visits]\label{lem:mot-delete}
\notformalized
For every labelling $\beta$, every set $R\subseteq[0,c_0]$ and every level $p$,
$\kappa(p)\le\kappa^R(p)+2\,\#(\Lambda(p)\setminus R)$. If $R$ is Borel, then
\[
  \int\kappa\,\mathrm dp\;\le\;\int\kappa^R\,\mathrm dp+2\,\nu([0,c_0]\setminus R).
\]
\end{lemma}

\begin{informalproof}
If $\Lambda(p)$ is infinite, then $\kappa(p)=0$. Otherwise every consecutive pair counted by $\kappa(p)$ but not by
$\kappa^R(p)$ has a point in $\Lambda(p)\setminus R$, and each point of $\Lambda(p)$ lies in at most two consecutive
pairs. For Borel $R$ the function $p\mapsto\#(\Lambda(p)\setminus R)$ is measurable with integral
$\nu([0,c_0]\setminus R)$, by \cref{lem:mot-levels}(a), and the lower integral of a sum is the sum of the lower integrals
when one summand, here $2\,\#(\Lambda(p)\setminus R)$, is measurable.
\end{informalproof}

The pointwise inequality is \leanname{TransposeFollowup.AffineMotion.levelCross_delete_le} in the second paper's Lean
library. The next lemma is the finite step of that paper's proof of its affine-time obstruction. For finitely many affine
pieces, the visits of a level in two pieces on which the path is not constant are related by an affine map of time; the
finitely many levels of the constant pieces form a null set of levels. The labels are built from F\o lner sets for the
finitely many affine maps that arise; such sets exist because the affine group is solvable, and the second paper
constructs them directly.

\begin{lemma}[Finite affine labels]\label{lem:mot-finite-labels}
\cited{the finite step of the proof of~\cite[Theorem~6.2]{SecondPaper2026}, which applies it to open pieces.
Lean-verified there for measurable pieces as
\leanname{TransposeFollowup.AffineMotion.exists_finite_affine_branch_small_cut}; the countable step of that proof,
\leanname{TransposeFollowup.AffineMotion.exists_countable_affine_path_cut_bound}, uses it. Not reproved here.}
Let $(P_t)_{t<k}$ be arbitrary functions on $[0,c_0]$, with $\Lambda_t(p)$, consecutive pairs and $\kappa_t^{R}$ defined
as above (no continuity is needed), and for each $t$ let $J_{t,1},\dots,J_{t,n}$ be Borel subsets of
$[0,c_0]$ on each of which $P_t$ agrees with an affine function; put $R_t=J_{t,1}\cup\dots\cup J_{t,n}$. For all
$\theta>0$ and $\varepsilon>0$ there is a labelling $\beta$ whose classes have measure at most $\theta$ and
$\sum_{t<k}\int\kappa_t^{R_t}\,\mathrm dp<\varepsilon$.
\end{lemma}

In the Lean statement the pieces need only be measurable, they may overlap and may carry constant functions, and the
$P_t$ need not be continuous. At each level it counts the crossing pairs with both points retained, and this count is at
least $\kappa_t^{R_t}(p)$.

The second form of the labels concerns first returns. For $u\in[0,c_0]$ let
$\psi(u)=\inf\{v\in[0,c_0]:v>u,\ P(v)=P(u)\}$, the \emph{first return} of the path to the level of $u$, defined on the
set $\operatorname{dom}\psi$ of those $u$ for which this set is nonempty, and let
$D^+=\{u\in\operatorname{dom}\psi:\psi(u)>u\}$. For a labelling $\beta$, the \emph{crossing set}
$A(\beta)=\{u\in D^+:\beta(\psi(u))\ne\beta(u)\}$ consists of the times whose first return changes the label. For a
family $(P_t)$ we write $\psi_t$, $D^+_t$ and $A_t(\beta)$.

\begin{lemma}[Returns]\label{lem:mot-returns}
\notformalized
\begin{enumerate}[label=(\alph*)]
\item $\operatorname{dom}\psi$ and $D^+$ are Borel sets, and $\psi\colon\operatorname{dom}\psi\to[0,c_0]$ is a Borel
map.
\item $P(\psi(u))=P(u)$ for every $u\in\operatorname{dom}\psi$, and $\psi$ is injective on $D^+$.
\item $\nu(\psi^{-1}(W)\cap D^+)\le\nu(W)$ for every Borel set $W\subseteq[0,c_0]$.
\item If $\Lambda(p)$ is finite, its consecutive pairs are exactly the pairs $(s,\psi(s))$ with $s\in\Lambda(p)\cap D^+$.
\item For every labelling $\beta$ the crossing set $A(\beta)$ is Borel, and $\int\kappa\,\mathrm dp\le\nu(A(\beta))$.
\end{enumerate}
\end{lemma}

\begin{informalproof}
(a) For $a\in\R$, the set $\{u\in\operatorname{dom}\psi:\psi(u)<a\}$ is the projection to the first coordinate of
$K_a=\{(u,v)\in[0,c_0]^2:u<v<a,\ P(u)=P(v)\}$: if $\psi(u)<a$, some $v$ in the set whose infimum is $\psi(u)$ is below
$a$, and conversely. $K_a$ is the intersection of a closed set with an open subset of the compact square, hence
$\sigma$-compact, so its projection is $\sigma$-compact and therefore Borel. Hence
$\operatorname{dom}\psi=\{\psi<c_0+1\}$ is Borel and $\psi$ is Borel on it, and so is
$D^+=\{u\in\operatorname{dom}\psi:\psi(u)>u\}$.
(b) $\psi(u)$ is the infimum of a nonempty subset of the closed set $\Lambda(P(u))\cap[u,c_0]$, so it lies in that set.
Let $u<u'$ be points of $D^+$ with $\psi(u)=\psi(u')=v$. Then $P(u)=P(v)=P(u')$, so $u'$ belongs to the set whose
infimum is $\psi(u)$, and $\psi(u)\le u'<\psi(u')=v$, a contradiction.
(c) The set $\psi^{-1}(W)\cap D^+$ is Borel by (a), and by (b) $\psi$ maps $\psi^{-1}(W)\cap D^+\cap\Lambda(p)$
injectively into $W\cap\Lambda(p)$ for every $p$. Integrate the inequality
$\#(\psi^{-1}(W)\cap D^+\cap\Lambda(p))\le\#(W\cap\Lambda(p))$ over $p$ and apply \cref{lem:mot-levels}(a) to both sides.
(d) If $\Lambda(p)$ is finite and $s\in\Lambda(p)$ is not its largest point, then $\psi(s)$ is the next point of
$\Lambda(p)$, so $\psi(s)>s$; the largest point is not in $\operatorname{dom}\psi$.
(e) With $Z$ as in the definition of a labelling, $\beta$ is constant on each of the finitely many components of
$[0,c_0]\setminus Z$, so it is Borel on $[0,c_0]$, and $A(\beta)$ is the union over $\ell$ of the sets
$\{u\in D^+:\beta(u)=\ell,\ \beta(\psi(u))\ne\ell\}$, which are Borel by (a). By (d),
$\kappa(p)=\#(A(\beta)\cap\Lambda(p))$ whenever $\Lambda(p)$ is finite, and $\kappa(p)=0$ otherwise. So
$\int\kappa\,\mathrm dp\le\nu(A(\beta))$, by monotonicity of the lower integral and \cref{lem:mot-levels}(a).
\end{informalproof}

\begin{lemma}[Finite affine labels, measure form]\label{lem:mot-finite-measure}
\cited{the finite step of the proof of~\cite[Theorem~6.2]{SecondPaper2026}, which applies it to push-forwards of
height measure; Lean-verified there, for every measure $\mu$ as in the statement, as
\leanname{TransposeFollowup.AffineMotion.exists_finite_affine_small_cut}. Not reproved here.}
Let $\Gamma_0$ be a finite set of affine bijections of $\R$, let $K\ge0$, and let $\mu$ be a finite Borel measure on
$[0,c_0]^2$ that is concentrated on the union of the graphs $\{(x,g(x))\}$ of the maps $g\in\Gamma_0$ and satisfies
$\mu(E\times[0,c_0])\le K\lambda(E)$ for every Borel set $E$. For all $\theta>0$ and $\varepsilon>0$ there is a
labelling $\beta$ whose classes have measure at most $\theta$ and $\mu(\{(x,y):\beta(x)\ne\beta(y)\})<\varepsilon$.
\end{lemma}

In the second paper's Lean library, \cref{lem:mot-finite-labels} is derived from this lemma, by pushing the measure of
levels forward to pairs of times in two affine pieces. The two lemmas are the only places below where amenability
enters.

\subsection{Theorems C, C\texorpdfstring{$'$}{'} and Cq}\label{sec:mot-statements}

We first define the class of Theorem~C. Its members are home-based in the sense of \cref{sec:prelim-walks}, up to a
tolerance, and every excursion from the home is made of boundedly many runs at constant rate.

\begin{definition}[Home-based heads with affine runs]\label{def:mot-hbar}
Let $K\ge1$ be an integer and $\rho\ge0$. A function $x\colon\{0,\dots,T\}\to\R$, for instance a walk of length $T+1$, is
\emph{HB-AR$(K,\rho)$} if there are a point $h$, its \emph{home}, and times $0=a_0<a_1<\dots<a_m=T$ such that every
\emph{segment} $[a_{l-1},a_l]$ satisfies:
\begin{enumerate}[label=(\roman*)]
\item (home-based) $\abs{x(a_{l-1})-h}\le\rho$ or $\abs{x(a_l)-h}\le\rho$;
\item (affine runs) there are times $a_{l-1}=s_0\le s_1\le\dots\le s_K=a_l$ and rates $q_1,\dots,q_K\in[-1,1]$ such
that $\abs{x(s)-x(s_{j-1})-q_j(s-s_{j-1})}\le\rho$ whenever $s_{j-1}\le s\le s_j$.
\end{enumerate}
The intervals $[s_{j-1},s_j]$ are the \emph{runs} of the segment. A machine \emph{has HB-AR$(K,\rho)$ work heads at
size $d$} if on every matrix input of size $d$ the trajectory of every work head, up to the halting time, is
HB-AR$(K,\rho)$.
\end{definition}

\begin{remark}[What the class contains]\label{rem:mot-hbar}
A head that stays within distance $\rho/2$ of its home during a time interval contributes a segment with a single run
of rate~$0$, so stays at home cost nothing. A home-based walk in the sense of \cref{sec:prelim-walks}, with home $h$,
whose legs consist of at most $K$ runs at constant rate up to deviation $\rho$ is HB-AR$(K,\rho+2)$: cut the time axis
at the first step of every leg; each resulting segment is contained in the corresponding leg, or exceeds it by one step
at its end. If $x$ is HB-AR$(K,\rho)$ and $\abs{x-x'}\le\rho'$, then $x'$ is HB-AR$(K,\rho+2\rho')$. With $\rho(d)=o(L(d))$ the
class therefore contains the usual home-based schedules: sweeps out and back from a home at the end of the used tape,
with any number of passes; copying excursions with small zigzags; excursions of all sizes in any order; any coupling
between the heads; excursions on both sides of the home; and homes that depend on the size and on the input. Nothing is
assumed about the input head or the output. The class excludes curved macroscopic motion: a macroscopic pause inside an
excursion costs a run, and an excursion whose speed varies along its length, for instance unit moves alternating with
pauses at a slowly varying duty cycle, needs unboundedly many runs.
\end{remark}

For a machine $M$ let $T_M(d)$ be the largest halting time of $M$ on matrix inputs of size $d$ (infinite if some input
does not halt).

\begin{theorem}[Theorem C: home-based heads with affine runs]\label{thm:C}
\notformalized
Let $K\ge1$ and $\rho(d)=o(L(d))$, and let $M$ be a machine that transposes. Suppose that for all large $d$ the machine
$M$ has HB-AR$(K,\rho(d))$ work heads at size $d$. Then $T_M(d)/L(d)\to\infty$. In particular $M$ does not witness~(T).
\end{theorem}

The hypothesis concerns every work head of $M$; for an oblivious machine it is a condition on its single head
schedule. For an adaptive machine it must hold on every input, with $K$ and $\rho$ uniform; the proof uses it only on
the representatives of the popular trace classes.

\paragraph{Why it matters.} Theorem~C is an unconditional impossibility result for a natural class: it needs neither
depth reduction nor the crux. The class allows any number of passes, excursions at all scales in any interleaving, any
coupling of the heads and arbitrary motion of the input head, and it includes adaptive machines. For oblivious machines whose work heads are all
home-based, \cref{cor:B} reaches the same conclusion for every motion, but only if the crux (\cref{def:crux}) holds;
Theorem~C needs no coding statement, for the subclass with affine runs. For a construction it
says that the work heads cannot all be home-based with legs made of boundedly many constant-rate runs. By Theorem~Cq,
in every limit, motion that is not countably-piecewise affine must carry at least half a block unit of travel in
total; for a home-based head such motion needs legs with unboundedly many runs whose speed varies along the leg, as
produced by counters or rulers.

\begin{theorem}[Theorem C$'$: open affine pieces]\label{thm:mot-Cprime}
\notformalized
Let $M$ be a machine that transposes. Suppose that for every $c>0$, every limit $(P_t)_{t<k}$ of $M$ at time constant
$c$ admits relatively open sets $U_t\subseteq[0,c_0]$ such that $P_t$ is affine on each connected component of $U_t$ and
$\nu_t([0,c_0]\setminus U_t)=0$. Then $T_M(d)/L(d)\to\infty$.
\end{theorem}

Theorem~C$'$ allows the non-affine set $[0,c_0]\setminus U_t$ to have positive Lebesgue measure, provided the head is
stationary on almost all of it. Theorem~C follows from the quantitative statement below, and so does C$'$.

\paragraph{Why it matters.} Theorem~C$'$ is the form of Theorem~Cq that is easiest to check on a concrete schedule: a
construction whose limit motion is affine on open intervals wherever the heads move fails, even when the set on which
the heads rest is a fat Cantor set (\cref{rem:mot-fatcantor}).

\begin{theorem}[Theorem Cq: curved motion]\label{thm:Cq}
\notformalized
Let $(P_t)_{t<k}$ be $1$-Lipschitz functions on $[0,c_0]$ that satisfy $(\star)_\emptyset$; by \cref{thm:Mstar} this
holds for every limit of every machine. For each $t$ let $E_{t,0},E_{t,1},\dots$ be Borel subsets of $[0,c_0]$ on each
of which $P_t$ agrees with an affine function, and let $N_t=[0,c_0]\setminus\bigcup_aE_{t,a}$. Then
\[
  \sum_{t<k}\nu_t(N_t)\;\ge\;\tfrac12 .
\]
\end{theorem}

In words: in every limit of a linear-time transposer, the limit paths of the work heads travel a total distance of at
least $\frac12$, in block units, while moving in a way that is not countably-piecewise affine. Since distance $1$
corresponds to about $L/G$ cells, this is about $L/(2G)$ cells of curved travel.

\paragraph{Why it matters.} Theorem~Cq is the quantitative statement behind Theorems~C and~C$'$, and it contains the
Lean-verified affine-time obstruction of the second paper (\cref{cor:mot-lambda}). It tells a construction how much
curved motion it needs: at least half a block unit of travel, whatever else the machine does. It does not exclude
curved motion (\cref{rem:mot-scope}).

\subsection{Proofs of Theorems Cq, C\texorpdfstring{$'$}{'} and C}\label{sec:mot-proofs}

We prove Theorem~Cq (\cref{thm:Cq}) first. The proof is the second paper's truncation argument with the speed measure in
place of Lebesgue measure: finitely many pieces are handled by \cref{lem:mot-finite-labels}, and the returns that involve
the other times are charged by their speed.

\begin{informalproof}
Fix $\theta\in(0,\frac12)$, $\varepsilon>0$ and $n\ge1$, and let $R_{t,n}=E_{t,0}\cup\dots\cup E_{t,n-1}$.
\Cref{lem:mot-finite-labels}, applied to these pieces, gives a labelling $\beta$ whose classes have measure at most
$\theta$ and $\sum_t\int\kappa_t^{R_{t,n}}\,\mathrm dp<\varepsilon$. By \cref{lem:mot-delete},
$\int\kappa_t\,\mathrm dp\le\int\kappa_t^{R_{t,n}}\,\mathrm dp+2\nu_t([0,c_0]\setminus R_{t,n})$ for every $t$, so
$(\star)_\emptyset$, applied to $\beta$, gives
\[
  1\;\le\;\varepsilon+2\theta+2\sum_{t<k}\nu_t\bigl([0,c_0]\setminus R_{t,n}\bigr).
\]
For each $t$ the sets $[0,c_0]\setminus R_{t,n}$ decrease to $N_t$ as $n\to\infty$, and $\nu_t$ is a finite measure, so
$\nu_t([0,c_0]\setminus R_{t,n})\to\nu_t(N_t)$. Letting $n\to\infty$ and then $\theta,\varepsilon\to0$ gives
$\sum_t\nu_t(N_t)\ge\frac12$.
\end{informalproof}

Theorem~C$'$ (\cref{thm:mot-Cprime}) follows from Theorem~Cq and the existence of limits.

\begin{informalproof}
Suppose not. Then there are $c>0$ and sizes $d_i\to\infty$ with $T_M(d_i)\le cL(d_i)$, so $M$ transposes and halts
within $cL(d_i)$ steps on every matrix input of size $d_i$, and by \cref{thm:Mstar}(b) it has a limit $(P_t)_{t<k}$ at
time constant $c$. Let $U_t$ be as in the hypothesis. The components of $U_t$ are countably many relatively open
intervals, on each of which $P_t$ is affine. By \cref{thm:Cq}, applied to these intervals, $\sum_t\nu_t([0,c_0]\setminus
U_t)\ge\frac12$, which contradicts $\nu_t([0,c_0]\setminus U_t)=0$.
\end{informalproof}

The proof of Theorem~C (\cref{thm:C}) rests on the structure of the limits of HB-AR heads: they sit at a home level,
apart from countably many \emph{teeth}, each piecewise affine with boundedly many pieces. The teeth may be dense in time.

\begin{lemma}[Limits of home-based heads with affine runs]\label{lem:mot-teeth}
\notformalized
Let $d_i\to\infty$, $c>0$, $K\ge1$ and $\rho_i=o(L(d_i))$, and let $x_i\colon\{0,\dots,T_i\}\to\Z$ be walks with
$x_i(0)=0$ and $1\le T_i\le cL(d_i)$ that are HB-AR$(K,\rho_i)$. Extend $x_i$ by $x_i(s)=x_i(T_i)$ for $s>T_i$, and suppose
that $P_i(u)=d_i^{-1}\lfloor x_i(\lfloor ud_i\rfloor r_{d_i})/r_{d_i}\rfloor$ converges uniformly on $[0,c_0]$ to a
function $P$. Then there is a level $p\in\R$ with the following properties.
\begin{enumerate}[label=(\alph*)]
\item The set $\{u\in[0,c_0]:P(u)\ne p\}$ is the union of countably many disjoint relatively open intervals of
$[0,c_0]$, the \emph{teeth}.
\item For each tooth $I$ there is a set $\Phi_I\subseteq I$ of at most $2K$ points such that $P$ is affine on each
component of $I\setminus\Phi_I$.
\item For every $y>0$, at most $c_0/(2y)+2$ teeth contain a point $u$ with $\abs{P(u)-p}>y$.
\item $\nu(\{u\in[0,c_0]:P(u)=p\})=0$, where $\nu=\abs{P'}\,\mathrm du$.
\end{enumerate}
\end{lemma}

\begin{informalproof}
Write $r_i=r_{d_i}$ and $R_i=R_{d_i}$. Let $X_i\colon[0,\infty)\to\R$ interpolate $x_i$ linearly between integers, and
let $Q_i(u)=X_i(uR_i)/R_i$; it is $1$-Lipschitz. If $s=\lfloor ud_i\rfloor r_i$, then $0\le uR_i-s<r_i$, so
$\abs{X_i(uR_i)-x_i(s)}\le r_i$, and $\abs{P_i(u)-x_i(s)/R_i}\le1/d_i$. Hence $\abs{P_i-Q_i}\le2/d_i$, so $Q_i\to P$
uniformly on $[0,c_0]$, and $P$ is $1$-Lipschitz (as in~\cite[Lemma~7.1(a)]{FirstPaper2026}).

\emph{Rescaled segments.} Let $h_i$ and $0=a_0<\dots<a_m=T_i$ be a home and segments for $x_i$ as in
\cref{def:mot-hbar}, and put $p_i=h_i/R_i$ and $\varepsilon_i=\rho_i/R_i$; then $\varepsilon_i\le G\rho_i/L(d_i)\to0$.
Divided by $R_i$, the segments become closed intervals that partition $[0,T_i/R_i]$, and we add the interval
$[T_i/R_i,c_0]$, on which $Q_i$ is constant, to the last segment as one more run, of rate~$0$; recall
$T_i/R_i\le cG<c_0$. On a run $[s_{j-1},s_j]$ the function $X_i(s)-X_i(s_{j-1})-q_j(s-s_{j-1})$ is linear between
consecutive integers, so its absolute value is at most $\rho_i$ at all real $s$ of the run. We obtain a partition of
$[0,c_0]$ into closed intervals, the segments, such that each segment has an endpoint at which
$\abs{Q_i-p_i}\le\varepsilon_i$ (for the extended last segment, recall that $Q_i$ is constant on the added interval),
and each segment is the union of at most $K$ consecutive closed intervals, its runs, on each of which
$\abs{Q_i(u)-Q_i(u')-q(u-u')}\le2\varepsilon_i$ for some $q\in[-1,1]$; the last segment has at most $K+1$ runs. The
first segment has an endpoint $s$ with $\abs{x_i(s)-h_i}\le\rho_i$, and $\abs{x_i(s)}\le s\le T_i$, so
$\abs{p_i}\le cG+\varepsilon_i$. Passing to a subsequence, which does not change $P$, we may assume that $p_i\to p$.

(a) The set $\{P\ne p\}$ is relatively open in $[0,c_0]$; its components are disjoint relatively open intervals, and
there are countably many of them.

(b) Let $I$ be a tooth and $[\alpha,\beta]\subseteq I$ a compact interval with $\alpha<\beta$. Then
$\delta=\min_{[\alpha,\beta]}\abs{P-p}>0$, and for large $i$ we have $\sup_{[0,c_0]}\abs{Q_i-P}<\delta/4$,
$\abs{p_i-p}<\delta/4$ and $\varepsilon_i<\delta/2$, so that $\abs{Q_i-p_i}>\delta/2>\varepsilon_i$ on $[\alpha,\beta]$.
Let $J_i$ be the component of $\{u\in[0,c_0]:\abs{Q_i(u)-p_i}>\varepsilon_i\}$ that contains $[\alpha,\beta]$. At most
two segments meet $J_i$: if a segment lies strictly between two segments that meet $J_i$, then both of its endpoints lie
between points of $J_i$, hence in the interval $J_i$, against the home-based condition. So $[\alpha,\beta]$ is covered
by the at most $2K+1$ runs of these two segments, and there are points
$\alpha=\tau_{i,0}\le\tau_{i,1}\le\dots\le\tau_{i,2K+1}=\beta$ and slopes $q_{i,j}\in[-1,1]$ such that
$\abs{Q_i(u)-Q_i(u')-q_{i,j}(u-u')}\le2\varepsilon_i$ for $u,u'\in[\tau_{i,j-1},\tau_{i,j}]$. Along a further
subsequence $\tau_{i,j}\to\tau_j$ and $q_{i,j}\to q_j$. If $\tau_{j-1}<u<u'<\tau_j$, then $u,u'\in[\tau_{i,j-1},\tau_{i,j}]$
for large $i$, and letting $i\to\infty$ gives $P(u')-P(u)=q_j(u'-u)$. So $P$ is affine near every point of
$(\alpha,\beta)\setminus\{\tau_1,\dots,\tau_{2K}\}$. Let $\Phi_I$ be the set of points $u$ in the interior of $I$ such
that $P$ is affine on no neighbourhood of $u$. Every compact interval in $I$ contains at most $2K$ points of $\Phi_I$ in
its interior, so $\Phi_I$ has at most $2K$ points. On the interior of each component of $I\setminus\Phi_I$ the function
$P$ is locally affine, hence affine, and by continuity it is affine on the whole component.

(c) A tooth that does not contain $0$ or $c_0$ has its endpoints at level $p$. If it contains a point $u$ with
$\abs{P(u)-p}>y$, then, $P$ being $1$-Lipschitz, $u$ is at distance more than $y$ from both endpoints, so the tooth has
length more than $2y$. The teeth are disjoint, and at most two of them contain $0$ or $c_0$.

(d) This is \cref{lem:mot-levels}(b).
\end{informalproof}

Theorem~C now follows from the lemma and Theorem~Cq.

\begin{informalproof}
Suppose not. Then there are $c>0$ and sizes $d_i\to\infty$ with $T_M(d_i)\le cL(d_i)$, and by \cref{thm:Mstar}(b) the
machine $M$ has a limit $(P_t)_{t<k}$ at time constant $c$ along a subsequence, with representatives $w_i$. For large
$i$ and each $t$, the trajectory $x_{t,w_i}$ up to the halting time $T_i\le cL(d_i)$ starts at cell $0$ and is
HB-AR$(K,\rho(d_i))$, and $T_i\ge L(d_i)\ge1$ because $M$ writes $L(d_i)$ output bits, at most one per step. So
\cref{lem:mot-teeth} applies to these trajectories. Hence there is a level $p_t$ such that the components of
the sets $I\setminus\Phi_I$, for the teeth $I$ of $P_t$, form countably many relatively open intervals on each of which
$P_t$ is affine, and they cover $[0,c_0]$ except for the set $\{P_t=p_t\}$ and countably many points. Both exceptional
sets are $\nu_t$-null, by \cref{lem:mot-teeth}(d) and \cref{lem:mot-levels}(c). So \cref{thm:Cq}, applied to these
intervals, gives $0\ge\frac12$, a contradiction.
\end{informalproof}

\subsection{Relation to the Lebesgue form}\label{sec:mot-lambda}

\begin{corollary}[Speed-weighted and Lebesgue forms]\label{cor:mot-lambda}
\notformalized
Let $(P_t)_{t<k}$ be $1$-Lipschitz functions on $[0,c_0]$ that satisfy $(\star)_\emptyset$, and let
$U_t\subseteq[0,c_0]$ be relatively open sets on which $P_t$ is locally affine. Then
$\sum_t\nu_t([0,c_0]\setminus U_t)\ge\frac12$, and hence $\sum_t\lambda([0,c_0]\setminus U_t)\ge\frac12$.
\end{corollary}

\begin{informalproof}
The components of $U_t$ are countably many relatively open intervals. On each of them $P_t$ is locally affine, hence
affine: its slope is locally constant on a connected set, hence constant, and then $P_t$ minus a linear function is
locally constant, hence constant. \Cref{thm:Cq} applied to these intervals gives the first inequality, and
$\nu_t\le\lambda$ gives the second.
\end{informalproof}

The second inequality is the quantitative affine-time obstruction of the second paper~\cite[Theorem~6.2]{SecondPaper2026},
which is Lean-verified there for paths satisfying its restricted star condition~\cite[Definition~6.1]{SecondPaper2026}.
That condition includes $(\star)_\emptyset$ as its case $S=\emptyset$ and is formally stronger; the corollary needs only
$(\star)_\emptyset$ and allows relatively open sets, so it slightly generalizes that theorem. Theorem~Cq carries the same
argument, F\o lner labels for finitely many affine pieces followed by truncation, from open sets and Lebesgue measure to
Borel pieces and the speed measure; the area formula charges the omitted returns by speed. The motion exclusions of~\cite[Sections~6 and~7]{SecondPaper2026}, for bounded words in explicit
operation palettes and for contracting recursions with regular stages, rest on that inequality.

The speed-weighted form is stronger exactly where the heads are slow on the non-affine set. For a limit of HB-AR heads,
\cref{lem:mot-teeth} shows that the largest relatively open set on which $P_t$ is locally affine misses only the
boundary of the home-level set $\{P_t=p_t\}$, a closed nowhere dense set, and countably many points. So the Lebesgue
form already excludes every machine with HB-AR work heads for which, in every limit, these boundaries have total
Lebesgue measure less than $\frac12$. Theorem~C is new for limits in which these nowhere dense sets have total measure at least $\frac12$, as
in the next example; they carry no speed, so for Theorem~Cq they cost nothing.

\begin{remark}[A fat Cantor schedule]\label{rem:mot-fatcantor}
Let $c$ satisfy $cG>1$, and let $C\subseteq[0,1]$ be a perfect nowhere dense set with $\lambda(C)=\frac12$ that
contains $0$ and $1$, for instance a Smith--Volterra--Cantor set. Let one work head do the following at size $d$: it stays at its home cell $0$
except during the time intervals $(aR_d,bR_d)$ for the gaps $(a,b)$ of $C$ of length at least $\delta_d$, where
$\delta_d\to0$ and $\delta_dR_d\to\infty$; during such an interval it moves away from home at unit speed until the
midpoint and back at unit speed. Every gap gives two segments of one run each, and the stays at home give segments with
one run of rate~$0$, so the trajectory is HB-AR$(1,1)$. The omitted gaps change the path by at most $\delta_d/2$, so the
rescaled trajectories converge uniformly to $P(u)=\operatorname{dist}(u,C)$ on $[0,1]$, constant afterwards. The
home-level set of $P$ within $[0,1]$ is $C$. No point of $C$ has a neighbourhood on which $P$ is affine: such an
affine function would vanish at that point and, $C$ being perfect, at points of $C$ arbitrarily close to it, hence
vanish identically, and $C$ would contain an interval. So every relatively open set on which $P$ is locally affine misses $C$, and the Lebesgue form gives nothing
because $\lambda(C)=\frac12$. But $\nu(C)=0$, so Theorem~C and Theorem~Cq exclude every machine all of whose work heads move
like this. Alternatively, $P$ is constant on the Borel set $C$, so $C$ can be taken as one of the pieces in Theorem~Cq;
the uncovered set is then countable, and even the Lebesgue-measure bound $\sum_t\lambda(N_t)\ge\frac12$ with Borel pieces,
which follows from Theorem~Cq because $\nu\le\lambda$, excludes this schedule.
\end{remark}

\begin{remark}[What escapes]\label{rem:mot-scope}
Theorem~Cq demands a minimum amount of speed-weighted motion that is not countably-piecewise affine; it says nothing
against curved motion. Theorem~Cr (\cref{thm:Cr}) excludes curved motion only when, in some clock, countably many
affine maps give the first returns at all times except a set carrying less than one block unit of travel.
A machine with one head whose home drifts, or whose motion is not home-based at all, falls outside Theorem~C even if
all its other heads are HB-AR; it is excluded only through Theorems~Cq and~Cr, that is, only if its limits are
countably-piecewise affine on enough of their speed, or have affine returns in a common clock on enough of it. Motion
that satisfies $(\star)$ exists, and by Theorems~Cq and~Cr it is necessarily curved, and in every clock its first
returns are not countably-piecewise affine on a set carrying at least one block unit of travel. The second paper constructs, for every machine, a fixed wrapper
that adds linear time and whose complete runs satisfy, whatever computation follows, the cut inequality
of~\cite[Definition~7.4]{FirstPaper2026} with zero slack, and also the conclusions of Theorem~M* when that computation
runs in linear time~\cite[Theorem~8.1]{SecondPaper2026}; this result is Lean-verified there. The heads of the wrapper
move along rational near-identity M\"obius maps (see~\cite[Section~13]{FirstPaper2026}). So Theorem~Cq is a constraint
on how a linear-time transposer moves its heads, not a lower bound by itself.
\end{remark}

\begin{remark}[How a machine can move along curves]\label{rem:mot-curved}
Proof sketch only: we indicate why curved macroscopic motion can be produced at amortised constant cost per step; this
is a design sketch, not a result. A bounded number of constant-rate runs per leg gives piecewise affine limits, so a
curved macroscopic leg needs unboundedly many runs whose speed, or duty cycle, varies along the leg. One way to produce
this is to hold two numbers $a,b$ of order $L^{\delta}$, for some
$0<\delta<1$, in unary on counter tapes and to alternate $b$ moves of one head with $a$ moves of another, each timed by
a sweep of a counter head; $a$ and $b$ change by single-cell edits scheduled by a slower accumulator. The two heads then
move at average speeds $b/(a+b)$ and $a/(a+b)$, with position error $o(L)$. For a M\"obius profile $f$ the function $1/f'$ is a quadratic
polynomial, so a two-level hierarchy of unary accumulators, a discrete second difference, suffices. This realizes
curved macroscopic motion only; the cell-exact pairings that data movement needs are a separate matter.
\end{remark}

\subsection{Consequences}\label{sec:mot-consequences}

The consequences concern the number of tapes that must interact, affine returns and changes of clock, and hierarchical
schedules; the first two, with Theorem~Cq, give the design constraint of \cref{cor:mot-design}.

\begin{corollary}[One work tape never expands]\label{cor:mot-one-tape}
\notformalized
\begin{enumerate}[label=(\alph*)]
\item Let $P$ be $1$-Lipschitz on $[0,c_0]$. For all $\theta>0$ and $\varepsilon>0$ there is a labelling $\beta$ whose
classes have measure at most $\theta$ and $\int\kappa\,\mathrm dp\le\varepsilon$.
\item Let $(P_t)_{t<k}$ be $1$-Lipschitz functions on $[0,c_0]$ that satisfy $(\star)_\emptyset$. Then at least two of
the $P_t$ are not monotone, and the non-monotone ones do not all have the same level sets: they are not all of the form
$h_t\circ P$ for one of them, $P$, and injective maps $h_t$.
\item Consequently, if $M$ has one work tape and transposes, then $T_M(d)/L(d)\to\infty$. In every limit of a machine
that transposes within $cL(d)$ steps on all inputs, at least two work-head limits are non-monotone, and they do not all
have the same level sets.
\end{enumerate}
\end{corollary}

\begin{informalproof}
(a) Let $\delta>0$ be a band width, to be chosen at the end. By \cref{lem:mot-levels}(b) the level set $\Lambda(p)$ is
finite for almost every $p$, so for almost every $\eta\in[0,1)$ all the levels $p_j=(j+\eta)\delta$, $j\in\Z$, have finite level sets; fix
such an $\eta$. The bands $T_j=\{u\in[0,c_0]:p_j\le P(u)<p_{j+1}\}$ partition $[0,c_0]$, only finitely many of them are
nonempty, and each is a finite union of intervals, because $P$ is continuous and takes each value $p_j$ only finitely
often. In a band with $\lambda(T_j)\le\theta$ make no cut, and put $m_j=0$. In each other band choose
$m_j=\lceil\lambda(T_j)/\theta\rceil-1\le\lambda(T_j)/\theta$ cut times, the first times at which the continuous function
$x\mapsto\lambda(T_j\cap[0,x])$ takes the values $\theta,2\theta,\dots,m_j\theta$. Let $\beta$ label $u\in T_j$ by a
natural number that encodes, injectively, the pair formed by $j$ and the number of cuts of $T_j$ that are at most $u$,
and label every $u\notin[0,c_0]$ by $0$. Then $\beta$ is constant on each of finitely many intervals that partition
$[0,c_0]$, so it is piecewise constant off $\emptyset$ (take $Z$ to be their endpoints). Its classes have measure at most
$\theta$: within $T_j$ they are the parts of $T_j$ before the first cut, between consecutive cuts, and after the last
cut, of measures $\theta,\dots,\theta$ and $\lambda(T_j)-m_j\theta\le\theta$.
Let $p\in[p_j,p_{j+1})$ have a finite level set, and let $s<s'$ be consecutive points of $\Lambda(p)$. Both lie in $T_j$,
so $\beta(s)\ne\beta(s')$ only if a cut of $T_j$ lies in $(s,s']$. These intervals are disjoint for different
consecutive pairs of $\Lambda(p)$, so $\kappa(p)\le m_j$. Hence $\int\kappa\,\mathrm dp\le\delta\sum_jm_j\le\delta
c_0/\theta$, which is $\varepsilon$ for $\delta=\varepsilon\theta/c_0$.

(b) If $P_t$ is monotone, each of its level sets is empty, a point or an interval, so $\kappa_t(p)=0$ for every labelling
and every $p$. Now suppose that the non-monotone $P_t$, if there are any, are all of the form $h_t\circ P$ for one of
them, $P$, and injective maps $h_t$. The interval $I=P([0,c_0])$ is nondegenerate, and $h_t$ is continuous on it: for
closed $F\subseteq\R$ the set $\{q\in I:h_t(q)\in F\}=P(P_t^{-1}(F))$ is compact. Being injective, $h_t$ is strictly
monotone on $I$. So $\Lambda_t(h_t(q))=\Lambda(q)$ for $q\in I$, the other level sets of $P_t$ are empty, and
$\kappa_t(h_t(q))=\kappa(q)$ for every labelling. Take $\theta=\frac14$. By uniform continuity there is $\delta>0$ such
that each of these $h_t$ maps every subinterval of $I$ of length at most $\delta$ onto an interval of length at most
$\theta/(4kc_0)$. Let $\beta$ be the labelling of (a) for $P$ with band width $\delta$, or any labelling with classes of
measure at most $\frac14$ if all $P_t$ are monotone. For each band, $\kappa_t\le m_j$ on the interval
$h_t([p_j,p_{j+1})\cap I)$, and $\kappa_t=0$ off the union of these intervals, so
$\int\kappa_t\,\mathrm dp\le\sum_jm_j\,\theta/(4kc_0)\le1/(4k)$. Hence $\sum_t\int\kappa_t\,\mathrm dp\le\frac14$, and
$(\star)_\emptyset$ gives $1\le\frac14+\frac12$, a contradiction.

(c) If $M$ has one work tape and transposes, and $T_M(d_i)\le cL(d_i)$ along some $d_i\to\infty$, then by
\cref{thm:Mstar}(b) $M$ has a limit, which satisfies $(\star)_\emptyset$ and consists of a single path, against (b). The
second statement is (b) for limits, which satisfy $(\star)_\emptyset$ by \cref{thm:Mstar}(a).
\end{informalproof}

\paragraph{Why it matters.} For one work tape the statement is much weaker than the classical superlinear lower bounds
in neighbouring models: $\Omega(L^{3/2}/(\log L)^{1/2})$ steps with one work tape and no output tape~\cite{DMS1991}, and
$\Theta(L^{5/4})$ for Boolean matrices with one work tape and a write-only output tape~\cite{DietzfelbingerMaass1993}.
Its interest is the reason: label time by thin height bands of the head's own position, and cut each band into pieces of
small measure. Consecutive visits of a level lie in one band and get different labels only across a cut, and each cut
separates at most one consecutive pair at each level; so crossings are rare, however curved the motion. The expansion that $(\star)$ demands must come from the interplay of
the returns of two or more heads moving along non-monotone paths with different level sets. Stationary and monotone heads
contribute nothing, and heads that all move in lockstep, even after a fixed change of scale on each tape, cannot supply
the expansion by themselves.

The next theorem replaces the affine pieces of Theorem~Cq by affine returns, seen through a common change of clock.

\begin{theorem}[Theorem Cr: affine returns]\label{thm:Cr}
\notformalized
Let $(P_t)_{t<k}$ be $1$-Lipschitz functions on $[0,c_0]$ that satisfy $(\star)_\emptyset$. Let $\varphi$ be a
homeomorphism of $[0,c_0]$ onto an interval $J$ such that $\varphi^{-1}$ maps $\lambda$-null sets to $\lambda$-null sets,
for instance the identity; we call such a $\varphi$ a \emph{clock}. Let $\Gamma$ be a countable set of affine bijections
of $\R$. For each $t$ let
\[
  X_t=\{u\in D^+_t:\varphi(\psi_t(u))\ne g(\varphi(u))\text{ for every }g\in\Gamma\}
\]
be the set of times whose first return is not given, in the clock $\varphi$, by a map in $\Gamma$; it is Borel by
\cref{lem:mot-returns}(a). Then
\[
  \sum_{t<k}\nu_t(X_t)\;\ge\;1 .
\]
\end{theorem}

In words: in every limit of a linear-time transposer, in every clock and for every countable set of affine maps, the
times at which a head's next return to its current level is not given by one of these maps carry at least one block
unit of travel in total.

\begin{informalproof}
Translating $J$, and conjugating the maps in $\Gamma$ by the same translation, we may assume that $J=[0,c_1]$. The
finite measure $A\mapsto\lambda(\varphi^{-1}(A))$ on $J$ vanishes on $\lambda$-null sets, so by the Radon--Nikodym
theorem it has a density $\rho$ with respect to $\lambda$, which we take Borel and finite everywhere, and for every
$\theta>0$ there is $\delta>0$ with $\lambda(\varphi^{-1}(A))\le\theta$ whenever $\lambda(A)\le\delta$. Let
$\Gamma_1\subseteq\Gamma_2\subseteq\cdots$ be finite sets with union $\Gamma$. For $n\ge1$ let $B_{t,n}$ be the set of
$u\in D^+_t$ with $\rho(\varphi(u))\le n$ and $\varphi(\psi_t(u))=g(\varphi(u))$ for some $g\in\Gamma_n$; it is Borel by
\cref{lem:mot-returns}(a).

Fix $\theta\in(0,\frac12)$, $\varepsilon>0$ and $n\ge1$, and choose $\delta$ for $\theta$. Push the restriction of
$\nu_t$ to $B_{t,n}$ forward along $u\mapsto(\varphi(u),\varphi(\psi_t(u)))$ and sum over $t$. This gives a finite Borel
measure $\mu$ on $J^2$, concentrated on the graphs of the maps in $\Gamma_n$, and by \cref{lem:mot-levels}(c), for
every Borel set $E\subseteq J$,
\[
  \mu(E\times J)=\sum_t\nu_t\bigl(B_{t,n}\cap\varphi^{-1}(E)\bigr)
  \le k\,\lambda\bigl(\varphi^{-1}(E\cap\{\rho\le n\})\bigr)=k\int_{E\cap\{\rho\le n\}}\rho\,\mathrm d\lambda
  \le kn\,\lambda(E).
\]
\Cref{lem:mot-finite-measure}, on $J$ and with $\delta$ in place of $\theta$, gives a labelling $\beta'$ whose classes
have measure at most $\delta$ and $\mu(\{(x,y):\beta'(x)\ne\beta'(y)\})<\varepsilon$. Let $\beta=\beta'\circ\varphi$ on
$[0,c_0]$ and $\beta=0$ elsewhere. Since $\varphi$ is a homeomorphism, $\beta$ is piecewise constant off $\emptyset$,
with the finite set $\varphi^{-1}(Z')$, where $Z'$ is that of $\beta'$. Its classes are the preimages under $\varphi$ of
those of $\beta'$, so by the choice of $\delta$ they have measure at most $\theta$. A point $u\in B_{t,n}$ lies in $A_t(\beta)$ exactly when
$\beta'(\varphi(u))\ne\beta'(\varphi(\psi_t(u)))$, so
$\sum_t\nu_t(A_t(\beta)\cap B_{t,n})=\mu(\{(x,y):\beta'(x)\ne\beta'(y)\})<\varepsilon$. By $(\star)_\emptyset$ and
\cref{lem:mot-returns}(e),
\[
  1\;\le\;\sum_t\int\kappa_t\,\mathrm dp+2\theta\;\le\;\sum_t\nu_t(A_t(\beta))+2\theta
  \;<\;\sum_t\nu_t(D^+_t\setminus B_{t,n})+\varepsilon+2\theta .
\]
As $n\to\infty$, the sets $B_{t,n}$ increase to $D^+_t\setminus X_t$, so $\nu_t(D^+_t\setminus B_{t,n})\to\nu_t(X_t)$,
and $1\le\varepsilon+2\theta+\sum_t\nu_t(X_t)$. Letting $\theta,\varepsilon\to0$ proves the theorem.
\end{informalproof}

For $\Gamma=\emptyset$ the theorem says $\sum_t\nu_t(D^+_t)\ge1$. By \cref{lem:mot-levels}(a) and~(b) and
\cref{lem:mot-returns}(d) the left side is the integrated level excess $\sum_t\int(\#\Lambda_t(p)-1)^+\,\mathrm dp$. The
second paper also bounds it below by $1$, for limits of linear-time runs that satisfy the first paper's cut inequality
with zero slack~\cite[Section~8]{SecondPaper2026}, and the bound follows directly from $(\star)_\emptyset$, since
$\kappa_t\le(\#\Lambda_t-1)^+$. The constant is $1$, not the $\frac12$ of Theorem~Cq, because each crossing is charged
once, to the earlier of its two visits, whereas deleting a visit can break two consecutive pairs.

A countable union of countable sets is countable, so the sets $X_t$ have a smallest version. If countable sets
$\Gamma_j$ make $\sum_t\nu_t(X_t)$ tend to its infimum over all countable $\Gamma$, their union $\Gamma^*$ attains it, and
then, for every countable $\Gamma$, $\nu_t$-almost every time in the set $X_t$ for $\Gamma^*$ also lies in the set $X_t$
for $\Gamma$. So in every clock the sets $X_t$ for $\Gamma^*$ carry at least one block unit of travel in total, and on no
part of them that carries travel are the first returns, seen in that clock, countably-piecewise affine.

The hypotheses on the clock matter. A Cantor function of a null Cantor set is continuous and nondecreasing but constant
on each gap, and almost every time lies in a gap; by \cref{lem:mot-returns}(c), $\nu_t$-almost no first return lands in
the Cantor set. With it in place of $\varphi$, countably many translations give the first return of $\nu_t$-almost every
time, so $\sum_t\nu_t(X_t)=0$ for every family of paths, and the conclusion fails.

Affine pieces, possibly after a change of scale on the tape, produce affine returns.

\begin{lemma}[Returns between affine pieces]\label{lem:mot-transitions}
\notformalized
Let $P$ be $1$-Lipschitz on $[0,c_0]$, let $\varphi\colon[0,c_0]\to\R$ be a Borel map, and let $h\colon\R\to\R$ be
injective. Let $E_a$ ($a\in\N$) be Borel subsets of $[0,c_0]$ and $f_a$ affine maps with $P=h\circ f_a\circ\varphi$ on
$E_a$, and let $N=[0,c_0]\setminus\bigcup_aE_a$. Let $\Gamma$ be the countable set of the maps $f_b^{-1}\circ f_a$ for
which neither $f_a$ nor $f_b$ is constant. Then the set of $u\in D^+$ with $\varphi(\psi(u))\ne g(\varphi(u))$ for every
$g\in\Gamma$, which is Borel by \cref{lem:mot-returns}(a), has $\nu$-measure at most $2\nu(N)$.
\end{lemma}

\begin{informalproof}
If $f_a$ is constant, then $P$ is constant on $E_a$, so $E_a$ lies in a level set and $\nu(E_a)=0$ by
\cref{lem:mot-levels}(b). Let $E^*$ be the union of the $E_a$ with $f_a$ not constant. The Borel set
$N^*=[0,c_0]\setminus E^*$ contains $N$ and is contained in $N$ together with the $\nu$-null sets $E_a$ with $f_a$
constant, so $\nu(N^*)=\nu(N)$. Let $u\in D^+$ with $u\in E_a$ and $\psi(u)\in E_b$, where neither $f_a$ nor $f_b$ is
constant. By \cref{lem:mot-returns}(b), $h(f_b(\varphi(\psi(u))))=P(\psi(u))=P(u)=h(f_a(\varphi(u)))$, and since $h$ is
injective, $\varphi(\psi(u))=f_b^{-1}(f_a(\varphi(u)))$. So the set in question is contained in
$N^*\cup(\psi^{-1}(N^*)\cap D^+)$, which has $\nu$-measure at most $2\nu(N)$ by \cref{lem:mot-returns}(c).
\end{informalproof}

\begin{corollary}[A common clock change and changes of scale do not help]\label{cor:mot-clock}
\notformalized
Let $(P_t)_{t<k}$ be $1$-Lipschitz functions on $[0,c_0]$ that satisfy $(\star)_\emptyset$. Let $\varphi$ be a
homeomorphism of $[0,c_0]$ onto an interval $J$ such that $\varphi^{-1}$ maps $\lambda$-null sets to $\lambda$-null sets,
for instance an increasing homeomorphism with absolutely continuous inverse, and for each $t$ let $h_t$ be an injective
map of $\R$ into itself, for instance a homeomorphism of $\R$. For each $t$ let $E_{t,a}$ ($a\in\N$) be Borel subsets
of $J$ on each of which $P_t\circ\varphi^{-1}=h_t\circ f$ for an affine function $f$, and let
$N_t=J\setminus\bigcup_aE_{t,a}$. Then $\sum_t\nu_t(\varphi^{-1}(N_t))\ge\frac12$.
\end{corollary}

\begin{informalproof}
For each $t$ apply \cref{lem:mot-transitions} to $P_t$, $\varphi$, $h_t$ and the Borel pieces $\varphi^{-1}(E_{t,a})$,
whose union misses exactly $\varphi^{-1}(N_t)$, and let $\Gamma$ be the union over $t$ of the countable sets of affine
maps that it produces. The sets $X_t$ of \cref{thm:Cr} for $\varphi$ and $\Gamma$ lie in the sets that the lemma bounds
for head $t$, because $\Gamma$ contains the maps produced for $t$. So $\nu_t(X_t)\le2\nu_t(\varphi^{-1}(N_t))$, and
\cref{thm:Cr} gives $1\le2\sum_t\nu_t(\varphi^{-1}(N_t))$.
\end{informalproof}

With $\varphi$ and all $h_t$ the identity this is Theorem~Cq, which therefore also follows from Theorem~Cr. The corollary
covers, for example, a global time profile that slows all heads down and speeds them up together, a shared duty cycle,
and heads whose speed depends in a fixed way on their position on the tape.

\paragraph{Why it matters.} Curvature that is common to all heads, or that comes from a fixed change of scale
on each tape, does not create expansion: if one change of clock and such changes of scale make every head
countably-piecewise affine, up to sets carrying little speed, then $(\star)$ fails. Non-affine motion is therefore
necessary but not sufficient. What is required, by Theorem~Cr, is that in no common clock are the first returns of the
heads given by countably many affine maps, except at times that carry less than one block unit of travel. Behind this
is amenability: the affine maps of $\R$ form a solvable group, and \cref{lem:mot-finite-measure} builds its labels from
F\o lner sets of finitely many of them. Head-dependent clocks, $P_t=f_t\circ\varphi_t$ with countably-piecewise affine
$f_t$ and different $\varphi_t$, are not covered; a M\"obius clock on one head and affine motion on another produce M\"obius transitions,
which is the regime of the monotone expanders of Bourgain and Yehudayoff~\cite{BourgainYehudayoff2013}.

Next we consider hierarchical schedules in which every move is a sweep at constant speed and every excursion returns
exactly to where it started. They may be arbitrarily deep and interleaved.

\begin{definition}[Stroke trees]\label{def:mot-stroke}
Let $K,K'\ge1$ be integers, $\gamma\in(0,1)$, $\eta_0>0$, $\rho\ge0$ and $R>0$, and let $x\colon\{0,\dots,T\}\to\Z$ be a
walk of length $T+1$. For integers $0\le a\le b\le T$ the \emph{time interval} $[a,b)$ is $\{a,\dots,b-1\}$; a finite family of time
intervals \emph{tiles} $[a,b)$ if its members are nonempty and pairwise disjoint and their union is $[a,b)$. A time
interval $[a,b)$ is a \emph{pause} if $x$ is constant on $\{a,\dots,b\}$. A \emph{stroke tree} for $x$ with parameters
$(K,K',\gamma,\eta_0)$, tolerance $\rho$ and scale $R$ is a finite rooted forest whose nodes are \emph{strokes} and
\emph{V's}, each with a time interval, such that:
\begin{enumerate}[label=(\roman*)]
\item the intervals of the roots tile $[0,T)$, and every root is a stroke;
\item every stroke $S$ has a direction $\sigma_S\in\{-1,1\}$, and its interval is tiled by at most $K$ \emph{runs}, any
number of pauses, and the intervals of its children, which are V's; a run is a time interval $[a,b)$ such that
$\sigma_Sx$ is nondecreasing on $\{a,\dots,b\}$ and, for some $q\in[0,1]$, $\abs{x(s)-x(a)-\sigma_Sq(s-a)}\le\rho$ for
$a\le s\le b$;
\item every V, with interval $[a,b)$, satisfies $x(a)=x(b)$, and its interval is tiled by any number of pauses and the
intervals of its children, which are at most $K'$ strokes;
\item every root either has extent at least $\eta_0R$ or has no runs, in which case it is a \emph{home stroke}; and every
stroke $C$ that is a child of a V of a stroke $S$ other than a home stroke satisfies
$\operatorname{ext}(C)\le\gamma\operatorname{ext}(S)$.
\end{enumerate}
Here the \emph{extent} $\operatorname{ext}(S)$ of a stroke $S$ is the sum of $\abs{x(b)-x(a)}$ over its runs $[a,b)$.
\end{definition}

Every stroke moves monotonically through at most $K$ runs, with any number of excursions (V's) inserted between runs and
at its ends; every excursion returns exactly to its starting cell and consists of at most $K'$ smaller strokes. A home
stroke models a fixed home from which excursions of any size start.

\begin{proposition}[Stroke trees are affine or stationary]\label{prop:mot-strokes}
\notformalized
Let $K,K'\ge1$, $\gamma\in(0,1)$, $\eta_0>0$ and $\rho(d)=o(L(d))$. Let $M$ be a machine and $t$ one of its work heads,
and suppose that for all large $d$, on every matrix input of size $d$, the trajectory of head $t$ up to the halting time
has a stroke tree with parameters $(K,K',\gamma,\eta_0)$, tolerance $\rho(d)$ and scale $R_d$. Let $(P_s)_{s<k}$ be a
limit of $M$ at some time constant, and let $U$ be the union of the open intervals in $(0,c_0)$ on which $P_t$ is affine.
Then $P_t$ is affine on each component of $U$, and $\nu_t([0,c_0]\setminus U)=0$. Consequently, if every work head of
$M$ satisfies the hypothesis and $M$ transposes, then $T_M(d)/L(d)\to\infty$.
\end{proposition}

\begin{informalproof}
Write $P=P_t$, let $d_i$ and $w_i$ be as in \cref{def:mot-limit}, let $x_i=x_{t,w_i}$, frozen after its halting time
$T_i$, let $R_i=R_{d_i}$, and let $Q_i(u)=X_i(uR_i)/R_i$, where $X_i$ interpolates $x_i$ linearly between integers. As
in the proof of \cref{lem:mot-teeth}, $Q_i\to P$ uniformly on $[0,c_0]$. The set $U$ is open, $P$ is affine on each of
its components because it is locally affine there, and $Z=[0,c_0]\setminus U$ is closed.

\emph{Claim: $\abs{P(v)-P(w)}\le\lambda([w,v]\cap U)$ for all $0\le w<v\le c_0$.} Fix $w<v$ and
$\eta\in(0,\min(\eta_0,c_0))$, and fix a large $i$; we drop the index $i$ where this causes no ambiguity. Call a stroke
\emph{big} if its extent is at least $\eta R$ and \emph{small} otherwise; home strokes have no runs and play no role. Let
$w''=\min(\lceil wR\rceil,T)$ and $v''=\min(\lfloor vR\rfloor,T)$. Since $X$ is $1$-Lipschitz and constant after $T$,
$\abs{Q(v)-Q(w)}\le(\abs{x(v'')-x(w'')}+2)/R$. For large $i$ either $w''<v''$, or $w''=v''=T$ and $x(v'')=x(w'')$; so
assume $w''<v''$, and let $I=[w'',v'')$.

Every time $s\in I$ has a \emph{chain}: the root containing $s$; then the run, pause or V of that root containing $s$;
then, if this is a V, the pause or stroke of that V containing $s$; and so on, down to a run or a pause. Call a V
\emph{inside} if its interval is contained in $I$. The \emph{tile} of $s$ is the interval of the first inside V on its
chain, or, if there is none, the run or pause at the end of its chain. Two tiles are disjoint or equal: nested intervals
of the forest belong to a node and a descendant, and a tile that is an inside V contains the tiles of none of its
descendants, since their chains pass through it. So the sets $H\cap I$, for the tiles $H$, tile $I$ into intervals, and
$x(v'')-x(w'')$ is the sum of the displacements of $x$ over them. Inside V's and pauses have displacement $0$. A tile
that is a run belongs to a stroke $S$ none of whose ancestor V's is inside; call such strokes \emph{exposed}. Since $x$ is
monotone with unit steps on a run, its displacement over a subinterval $[a',b')$ of a run $[a,b)$ is at most
$\min(\abs{x(b)-x(a)},\,b'-a')$. Hence
\[
  \abs{x(v'')-x(w'')}\le\sum_{S\text{ exposed}}c(S),\qquad
  c(S)\le\min\bigl(\operatorname{ext}(S),\ \#(I\cap\mathrm{runs}(S))\bigr),
\]
where $\mathrm{runs}(S)$ is the union of the runs of $S$.

\emph{Small exposed strokes.} An exposed stroke that is not a root has a parent V that meets $I$ but is not inside;
this V contains $w''-1$ and $w''$, or it contains $v''-1$ and $v''$. The V's of the first kind contain $w''$, so they lie
on the chain of $w''$, and they are its first $J$ V's for some $J$, say $V_1\supset V_2\supset\dots\supset V_J$, with
$V_j$ a child of the stroke $S_{j-1}$ of that chain and $S_j$ a child of $V_j$. For $j\ge2$ the stroke $S_{j-1}$ is not a
root, hence not a home stroke, and repeated use of (iv) gives $\operatorname{ext}(S_{j-1})\le\gamma^{j-2}
\operatorname{ext}(S_1)\le\gamma^{j-2}T$. So every child of $V_j$ has extent at most $\gamma^{j-1}T\le\gamma^{j-1}c_0R$,
because $T\le cL\le c_0R$. Each $V_j$ has at most $K'$ children, so the small exposed children of $V_1,\dots,V_J$
contribute at most
\[
  K'\eta R+K'R\sum_{j\ge1}\min(\eta,\gamma^jc_0)\;\le\;K'\eta R\Bigl(1+\log_{1/\gamma}\frac{c_0}{\eta}+\frac1{1-\gamma}\Bigr):
\]
at most $\log_{1/\gamma}(c_0/\eta)$ terms of the sum equal $\eta$, and the others form a geometric series with first
term below $\eta$. The V's of the second kind lie on the chain of $v''-1$, and their small children contribute at most
the same amount. Every small exposed stroke is a child of one of these V's, because the roots are big or home strokes.
So the small exposed strokes contribute at most $\delta(\eta)R$, where
$\delta(\eta)=2K'\eta(1+\log_{1/\gamma}(c_0/\eta)+1/(1-\gamma))$ tends to $0$ with $\eta$.

\emph{Big strokes.} Runs of different strokes are disjoint, and a big stroke spends at least $\eta R$ steps in its runs.
Since $T\le c_0R$, there are at most $c_0/\eta$ big strokes and at most $m=\lfloor Kc_0/\eta\rfloor$ runs of big
strokes. List these runs as $[a_{i,l},b_{i,l})$, $l\le m$, padding with empty intervals, and let
$J_{i,l}=[a_{i,l}/R_i,\,b_{i,l}/R_i]$. Since $[w''/R,v''/R]\subseteq[w,v]$, the big exposed strokes contribute at most
$\sum_l\#(I\cap[a_{i,l},b_{i,l}))\le R_i\sum_l\lambda([w,v]\cap J_{i,l})$. Altogether
\[
  \abs{Q_i(v)-Q_i(w)}\;\le\;\delta(\eta)+\sum_{l\le m}\lambda\bigl([w,v]\cap J_{i,l}\bigr)+\frac2{R_i}.
\]

\emph{Limit.} By (ii), as in the proof of \cref{lem:mot-teeth}, $\abs{Q_i(u)-Q_i(u')-q_{i,l}(u-u')}\le2\rho(d_i)/R_i$
for $u,u'\in J_{i,l}$, with slopes $q_{i,l}\in[-1,1]$. Along a subsequence the endpoints of the $J_{i,l}$ and the
slopes $q_{i,l}$ converge, to intervals $J_l\subseteq[0,c_0]$ and slopes $q_l$. The $J_{i,l}$ have pairwise disjoint
interiors, hence so do the $J_l$. On the interior of each $J_l$ the function $P$ is affine with slope $q_l$, as in the
proof of \cref{lem:mot-teeth}, so this interior lies in $U$. Moreover $\lambda([w,v]\cap J_{i,l})\to\lambda([w,v]\cap
J_l)$, and $\sum_l\lambda([w,v]\cap J_l)\le\lambda([w,v]\cap U)$ because the interiors of the $J_l$ are disjoint and
lie in $U$. Letting $i\to\infty$ gives $\abs{P(v)-P(w)}\le\delta(\eta)+\lambda([w,v]\cap U)$, and letting $\eta\to0$
proves the claim.

\emph{Conclusion.} At almost every point $u$ of $Z$ the function $P$ is differentiable and $u$ is a Lebesgue density point
of $Z$. For such $u$ and small $\tau>0$ the claim gives
$2\tau\abs{P'(u)}+o(\tau)=\abs{P(u+\tau)-P(u-\tau)}\le\lambda([u-\tau,u+\tau]\setminus Z)=o(\tau)$, so $P'(u)=0$. Hence $\nu_t(Z)=0$. The
last sentence of the proposition follows from \cref{thm:mot-Cprime}, with $U_s$ the set $U$ of head $s$.
\end{informalproof}

\paragraph{Why it matters.} Stroke trees model hierarchical copy-and-return schedules of unbounded depth, with roots of
macroscopic size or a fixed home. The proposition says that, however deep and however interleaved, such motion is
affine or stationary in the limit, and hence excluded when all work heads move this way. The non-affine sets
$[0,c_0]\setminus U$ are nowhere dense, since $P_t'=0$ almost everywhere on them. The Lebesgue form already excludes stroke trees whose non-affine sets have total
Lebesgue measure less than $\frac12$; the proposition is needed when that measure is at least $\frac12$, as for
Cantor-type arrangements of V's. A construction must break one of
the hypotheses, for instance with strokes made of unboundedly many runs of varying density, which are curved; V's that
drift; V's with unboundedly many strokes, or nested strokes whose extents do not contract geometrically; or roots that
are not macroscopic, concatenated with varying density.

\begin{corollary}[The motion a linear-time transposer needs]\label{cor:mot-design}
\notformalized
Let $M$ transpose within $cL(d)$ steps on every matrix input of every size $d$, for instance a machine that
witnesses~(T) with constant $c/2$, and let $(P_t)_{t<k}$ be a limit of $M$ at time constant $c$; such limits exist along
a subsequence of every sequence of sizes. Then:
\begin{enumerate}[label=(\alph*)]
\item (curved) for all Borel sets $E_{t,a}$ on each of which $P_t$ is affine, the uncovered sets $N_t$ satisfy
$\sum_t\nu_t(N_t)\ge\frac12$;
\item (head-dependent) this survives every common change of clock and every fixed change of scale on each tape: for
every homeomorphism $\varphi$ of $[0,c_0]$ onto an interval $J$ such that $\varphi^{-1}$ maps null sets to null sets and
every countable set $\Gamma$ of affine bijections of $\R$, the times $u\in D^+_t$ with
$\varphi(\psi_t(u))\notin\{g(\varphi(u)):g\in\Gamma\}$ carry travel at least $1$ in total, that is, the sum over $t$
of their $\nu_t$-measures is at least $1$; in particular, for all
injective maps $h_t$ of $\R$ and all Borel sets $E_{t,a}\subseteq J$ on each of which $P_t\circ\varphi^{-1}=h_t\circ f$
for an affine $f$, the uncovered sets $N_t\subseteq J$ satisfy $\sum_t\nu_t(\varphi^{-1}(N_t))\ge\frac12$;
\item (interacting tapes) at least two of the $P_t$ are not monotone, and these do not all have the same level sets;
in particular
$k\ge2$.
\end{enumerate}
\end{corollary}

\begin{informalproof}
Limits exist by \cref{thm:Mstar}(b), and they satisfy $(\star)_\emptyset$ by \cref{thm:Mstar}(a). Then (a) is
\cref{thm:Cq}, (b) is \cref{thm:Cr,cor:mot-clock} and (c) is \cref{cor:mot-one-tape}(b).
\end{informalproof}

\paragraph{Why it matters.} This is the motion a construction must use: curved, with curvature that differs between
heads, on at least two work tapes whose returns interact. The monotone expanders of Bourgain and
Yehudayoff~\cite{BourgainYehudayoff2013}, driven along near-identity M\"obius maps as in the architecture
of~\cite[Section~13]{FirstPaper2026}, are consistent with all three conditions. The corollary constrains only the
macroscopic limit of the motion. Structure at mesoscopic scales, such as nested excursions much shorter than $L$, is
invisible in the limit, and a curved macroscopic host can carry any mesoscopic structure; so the motion constraints say
nothing about depth reduction (\cref{sec:depth}), which concerns all scales at once. In this sense the results of this
section are complementary to the decomposition of \cref{sec:decomposition}: they exclude classes of macroscopic motion
without any coding hypothesis, while the decomposition applies to every motion of an oblivious machine but needs depth
reduction and the crux.
